\section{Boundary and conclusion}

Finite regular receipt lifts of connected graphs admit a complete algebraic
description. Edge receipts define a representation of the graph fundamental
group. Vertex gauge changes only its conjugacy class. The image subgroup
simultaneously determines reachability, component size, and connectedness:
\[
\text{gauge class}=[\rho_{\alpha}]_R,
\qquad
\text{component count}=[R:H_{\alpha}],
\qquad
\text{connected}\Longleftrightarrow H_{\alpha}=R.
\]

The one-cycle multiplication law is therefore not an isolated mechanism. It
is the rank-one shadow of a graph-wide holonomy-subgroup theorem. General
graphs permit several circuit receipts to cooperate, allowing connected lifts
with noncyclic receipt groups whenever the visible cycle rank is large enough
to generate them.

The theorem is algebraic, not physical. It does not derive a base/lift
projection, receipt group, edge cocycle, action magnitude, time variable,
energy law, force, or physical carrier. Those belong to the raw-graph binding
problem. The supplied certificate schema turns that problem into an explicit
finite sequence of gates.

The earned statement is:

\begin{quote}
The graph supplies the possible circuits. The receipt representation records
what every circuit carries home. Its image determines which lifted states can
belong to one connected history.
\end{quote}
